% SGA 3, Expose V, Loop 1 English reconstruction.
% Scope: Section 7 and Theorem 7.1 through the reduction to local
% quasi-sections, including the opening seam on component-local p. 25.
% Authority: Polo--Gille Exp5-13oct24.pdf, component-local pp. 24--25
% (combined-reader pp. 286--287), re-edited headers 274--275.
% Jacob Reinhold's CC BY 4.0 Markdown is comparison/drafting material only.

\section*{7. 固有平坦群胚的商}
\phantomsection\label{sec:sga3-V-7}
\addcontentsline{toc}{section}{7. 固有平坦群胚的商}

\medskip
\phantomsection\label{thm:sga3-V-7.1}%
\noindent\textbf{定理 7.1.}---%
\footnote{\label{note:sga3-V-34}编者注: 这里提及
S.~Keel 和 S.~Mori 的文章 \hyperref[bib:sga3-V-KM97]{[KM97]},
该文建立了如下定理. 设 \(X\) 是局部 Noether 底概形 \(S\)
上的有限型代数空间, \(j:R\to X\times_S X\) 是平坦群胚,
其稳定子 \(j^{-1}(\Delta_X)\) 在 \(X\) 上有限.
则存在一个代数空间, 它既是 \(X\) 关于 \(R\) 的几何商,
也是一致范畴商; 此外, 若 \(j\) 分离, 则此商分离.
特别地, 若平坦 \(S\)-群概形 \(G\) 固有地作用于 \(X\),
且稳定子有限 (即
\[
  G\times_S X\longrightarrow X\times_S X,\qquad
  (g,x)\longmapsto(x,g\cdot x)
\]
固有, 且对角线的稳定子在 \(X\) 上有限), 则存在几何商 \(X\to X/G\).
当 \(G\) 为约化 \(S\)-群概形时, 这是 J.~Koll\'ar
\hyperref[bib:sga3-V-Ko97]{[Ko97]} 的结果.}%
\emph{设 \(S\) 是局部 Noether 概形, 并设}
\SGAThreeDiagram[0.56\textwidth]
  {figures/exp5/p274_theorem71_groupoid.png}
  {\hyperref[thm:sga3-V-7.1]{定理~7.1} 中的群胚.}
  {fig:sga3-V-7.1-groupoid}
  {第 V 篇, 分篇第~24 页 / 合订本第~286 页 / 重版本页眉页码~274}
\emph{是一个 \((\mathbf{Sch}/S)\)-群胚,
使得 \(d_1\) 固有且平坦, \(X_0\) 在 \(S\) 上拟射影,}%
\footnote{\label{note:sga3-V-35}编者注: 关于 \(X_0\) 的此假设是必要的;
见定理~\hyperref[thm:sga3-V-4.1]{4.1} 的\hyperref[note:sga3-V-17]{编者注~(17)}.}
\emph{且以 \(d_0\) 和 \(d_1\) 为分量的态射
\[
  d:X_1\longrightarrow X_0\times_S X_0
\]
拟有限. 则:}
\begin{enumerate}
\item[(i)]\phantomsection\label{thm:sga3-V-7.1-i}%
  \emph{\((d_0,d_1)\) 在 \((\mathbf{Sch}/S)\) 中有余核 \((Y,p)\).
  此外, 每个这样的 \((Y,p)\) 都是所有赋环空间范畴中
  \((d_0,d_1)\) 的余核.}
\item[(ii)]\phantomsection\label{thm:sga3-V-7.1-ii}%
  \emph{态射 \(p\) 满射、开、固有且有限呈示,
  而 \(Y\to S\) 分离且有限呈示.}%
  \footnote{\label{note:sga3-V-36}编者注: TDTE~III,
  \SGARefLink{lem:sga3-V-6.1}{定理~6.1} 断言, 若 \(S\) 为 Noether 概形,
  则 \(Y\to S\) 拟射影. 编者未能看出如何从拟截面的局部存在性推出此点.}
\item[(iii)]\phantomsection\label{thm:sga3-V-7.1-iii}%
  \emph{以 \(d_0\) 和 \(d_1\) 为分量的态射
  \[
    X_1\longrightarrow X_0\times_Y X_0
  \]
  满射.}
\item[(iv)]\phantomsection\label{thm:sga3-V-7.1-iv}%
  \emph{若 \((d_0,d_1)\) 是等价对,
  则 \(X_1\to X_0\times_Y X_0\) 是同构, 且 \(p\) 忠实平坦.}%
  \footnote{\label{note:sga3-V-37}编者注: 随后的推论已明确写出;
  见 \hyperref[bib:sga3-V-Ray67a]{[Ray67a]}, 定理~1(iii),
  以及\hyperref[thm:sga3-V-7.1]{此定理}证明的末尾.
  还需提及, 对等价关系的情形,
  \hyperref[bib:sga3-V-AK80]{[AK80]}, 定理~2.9
  给出了基于 Hilbert 概形存在性的\hyperref[thm:sga3-V-7.1]{定理~7.1} 的另一个证明;
  亦见 \hyperref[bib:sga3-V-BLR90]{[BLR90]}, \S8.2, 定理~12.
  在此情形下还证明了 \(Y\to S\) 拟射影.}
  \emph{此外, \((Y,p)\) 是 fppf 拓扑的层范畴中 \((d_0,d_1)\) 的余核,
  并且对每个基变换 \(Y'\to Y\), 概形 \(Y'\) 是群胚
  \(X_\ast\times_Y Y'\) 的余核,
  后者由 \(X_\ast\) 作基变换 \(X_0\times_Y Y'\to X_0\) 得到.}
\end{enumerate}

\emph{特别地, 对每个基变换 \(S'\to S\), 概形
\[
  Y'=Y\times_S S'
\]
是 \(S'\)-群胚
\[
  X'_\ast=X_\ast\times_S S'.
\]
的余核. 因此在此情形下, \textquotedblleft{}商的构造与基变换可交换\textquotedblright.}

设 \((Y,p)\) 是所有赋环空间范畴中 \((d_0,d_1)\) 的余核.
引理~\hyperref[sga3:exp5:lemma:1.2]{1.2} 表明,
为证明\hyperref[thm:sga3-V-7.1-i]{(i) 项},
只须证明每个点 \(z\in X_0\) 都有饱和开邻域 \(U_z\),
使得若以 \(\widetilde d_0\) 和 \(\widetilde d_1\)
表示 \(d_0\) 和 \(d_1\) 在
\[
  d_0^{-1}(U_z)=d_1^{-1}(U_z),
\]
上的限制, 并以 \((Q,q)\) 表示赋环空间范畴中
\((\widetilde d_0,\widetilde d_1)\) 的余核,
则 \(Q\) 是概形, \(q\) 是概形态射.

因此由引理~\hyperref[lem:sga3-V-6.1-i]{6.1(i)},
只须证明每个点 \(z\in X_0\) 都有饱和开邻域 \(U_z\),
使由 \(X_\ast\) 在 \(U_z\) 上诱导的群胚具有拟截面.
甚至可以假设 \(z\) 在其位于 \(S\) 上的纤维中闭
(我们称 \(z\) \emph{相对于 \(S\) 闭}).%
\footnote{\label{note:sga3-V-38}编者注: 事实上,
若对每个相对于 \(S\) 闭的点 \(z\) 都已构造这样的开邻域 \(U_z\),
则这些 \(U_z\) 的并覆盖 \(X_0\):
其闭补集在 \(S\) 上的每条纤维都是没有闭点的 Noether 概形, 因而为空.}
\(U_z\) 的存在性由下面的引理得到.
% SGA3-V-V2-EOL-PRESERVATION

% Translation coverage: entire assigned component; adjoining components retain the original statement/proof seams.
% Source lines: 122; translated body lines before this audit comment: 113.
% Math preserved: 74 inline formulas, 6 display blocks, 0 equation environments; mathematical prose translated where present.
% Labels: 11/11; hyperrefs: 13/13; SGARefLinks: 1/1; inline targets: 0/0; footnotes: 5/5.
% Figure arguments preserved; all rendered captions and natural-language source locators translated.
% Unresolved TODOs: 0; compilation and visual inspection remain with the volume integrator.
